\documentclass{article}
\usepackage[UTF8]{ctex}
\usepackage{amsmath,amssymb,amsfonts}
\usepackage{enumitem}
\usepackage{geometry}
\geometry{a4paper, margin=1in}
% === 错题本排版优化（不参与题目标识） ===
\setlength{\parindent}{0pt}
\setlength{\parskip}{0.35em plus 0.10em minus 0.08em}
\setlength{\emergencystretch}{3em}
\setlist{itemsep=0.25em, topsep=0.40em, parsep=0pt}
\allowdisplaybreaks[2]
% === 错题本排版优化结束 ===
\begin{document}
\section*{中南大学2026年高等代数真题}
\section*{137．中南大学2026年高等代数试卷}

\subsection{题目1}
1．（16 分）已知 \(a_{1}, a_{2}, \cdots, a_{n}\) 是互异的整数，证明：

\[
h(x)=\left(x-a_{1}\right)^{2}\left(x-a_{2}\right)^{2} \cdots\left(x-a_{n}\right)^{2}+1
\]

在有理数域 \(\mathbb{Q}\) 上不可约。

\subsection{题目2}
2．（16 分）设 \(n \geq 2\) ，且 \(a_{i} \neq 0(i=1,2, \cdots, n)\) ，计算行列式

\[
|A|=\left|\begin{array}{cccc}
0 & a_{1}+a_{2} & \cdots & a_{1}+a_{n} \\
a_{2}+a_{1} & 0 & \cdots & a_{2}+a_{n} \\
\vdots & \vdots & & \vdots \\
a_{n}+a_{1} & a_{n}+a_{2} & \cdots & 0
\end{array}\right|
\]

\subsection{题目3}
3．（16 分）设 \(a_{1}, a_{2}, \cdots, a_{n}\) 为 \(n(n \geq 4)\) 个互异的实数，求下列 \(n\) 个未知量 \(x_{1}, x_{2}, \cdots, x_{n}, n-2\) 个方程的齐次线性方程组的通解

\[
\left\{\begin{array}{c}
x_{1}+x_{2}+\cdots+x_{n}=0 \\
a_{1} x_{1}+a_{2} x_{2}+\cdots+a_{n} x_{n}=0 \\
a_{1}^{2} x_{1}+a_{2}^{2} x_{2}+\cdots+a_{n}^{2} x_{n}=0 \\
\cdots \cdots \\
a_{1}^{n-3} x_{1}+a_{2}^{n-3} x_{2}+\cdots+a_{n}^{n-3} a_{n}=0
\end{array}\right.
\]

\subsection{题目4}
4．（16 分）已知整数 \(n \geq 3\) ，实数 \(x_{1}, x_{2}, \cdots, x_{n}\) 满足 \(x_{1}+x_{2}+\cdots+x_{n}=0, x_{1}^{2}+x_{2}^{2}+\cdots+x_{n}=1\) ，证明：

\[
x_{1} x_{2}+x_{2} x_{3}+\cdots+x_{n-1} x_{n}+x_{n} x_{1} \leq \cos \frac{2 \pi}{n} .
\]

\subsection{题目5}
5．（16 分）设 \(A\) 是 \(m \times n\) 复矩阵，\(B\) 和 \(C\) 分别为 \(n\) 级和 \(m\) 级复矩阵，方阵 \(B, C\) 没有公共的特征值，而且 \(A B=C A\) ，证明：\(A=O\) 。

\subsection{题目6}
6．（18 分）设 \(A, B\) 为 \(n\) 阶复方阵，\(f_{A}(\lambda), f_{B}(\lambda)\) 分别为 \(A, B\) 的特征多项式，证明：

\begin{itemize}
  \item[（1）] \(f_{A}(B)\) 可逆的充分必要条件为 \(\left(f_{A}(\lambda), f_{B}(\lambda)\right)=1\) 。
  \item[（2）] 矩阵方程 \(A X-X B=O\left(X \in \mathrm{C}^{\mathrm{n} \times \mathrm{n}}\right)\) 有非零解的充分必要条件为 \(\left(f_{A}(\lambda), f_{B}(\lambda)\right) \neq 1\) 。
\end{itemize}

\subsection{题目7}
7．（16 分）设 \(V\) 是复数域 C 上的 \(n\) 维线性空间，\(f_{i}: V \rightarrow \mathrm{C}(i=1,2)\) 是非零线性函数，且不存在常数 \(c\)，使得 \(f_{1}=c f_{2}\) ，证明：任意的 \(\alpha \in V\) 都可表示为 \(\alpha=\alpha_{1}+\alpha_{2}\) ，使得 \(f_{1}(\alpha)=f_{1}\left(\alpha_{2}\right), f_{2}(\alpha)=f_{2}\left(\alpha_{1}\right)\)－

\subsection{题目8}
8．（20 分）设 \(\sigma, \tau\) 是数域 \(P\) 上的线性空间 \(V\) 上的线性变换，且 \(\sigma^{2}=\sigma\) ．证明：

\begin{itemize}
  \item[（1）] \(\operatorname{Ker} \sigma=\{a-\sigma(a) \mid a \in V\}\) 。
  \item[（2）] 线性空间 \(V\) 是否可以表示为 \(\operatorname{Ker} \sigma\) 与 \(\operatorname{Im} \sigma\) 的直和？若成立，请给出证明；若不成立，请举出反例．
  \item[（3）] 请给出 \(\operatorname{Ker} \sigma\) 与 \(\operatorname{Im} \sigma\) 为 \(\tau\)－子空间的充分必要条件，并证明．
\end{itemize}

\subsection{题目9}
9．（16 分）已知 \(A\) 是一个 \(n\) 级正定矩阵，证明：存在唯一的正定矩阵 \(G\) ，满足 \(A=C^{2}\) ．
\end{document}
